\scanpage{49}
\mspage{76}
\margin{n°382}
\[
\begin{aligned}
(f\omega)\wedge\omega'&=f.(\omega\wedge\omega')=\omega\wedge f\omega',\\
i(fs)\omega&=f.i(s)\omega=i(s)(f\omega),\quad\text{etc.}
\end{aligned}
\]
6.8.10. On définit de même la \uline{puissance symétrique} k-\uline{ième} $S^k(E)$ d'un espace fibré vectoriel E ; ses fibres sont égales aux puissances symétriques k-ièmes des fibres de E. On laisse au lecteur le soin de formuler, pour les puissances symétriques des résultats analogues à ceux énoncés ci-dessus pour les puissances extérieures.

6.9 -- \uline{Changement de structure}.

6.9.1. Soit E un espace fibré vectoriel de classe $C^r$ ; c'est aussi un espace fibré vectoriel de classe $C^{r'}$ pour tout $r'\leq r$, ainsi qu'un espace fibré vectoriel au sens topologique ; de même, s'il est analytique sur une extension de K, il l'est aussi sur K.

6.9.2. Les changements de structure définis au n° 6.9.1 sont compatibles avec la formation des images réciproques, ainsi qu'avec le \qquad procédé de construction du n° 6.6.
\begin{center}-:-:-:-\end{center}
\mspage{48}
